% !TEX root = main.tex
\section{서론}
\subsection{제품정보의 신뢰성과 공급망 근거}
디지털 제품 여권(Digital Product Passport, DPP)은 제품의 특성·환경정보·이력을 제품과 연결하여 필요한 주체가 조회할 수 있도록 하는 체계다.
EU 배터리 규정은 2027년 2월 18일부터 시장에 출시되거나 사용 개시되는 전기차용 배터리, 경량 운송수단용 배터리, 용량 2 kWh 초과 산업용 배터리에 전자 배터리 여권을 요구한다.
여권의 환경정보 항목에는 탄소발자국과 재활용 원료 함량이 포함되며, 각 정보 의무의 적용은 해당 조항의 일정에 따른다. [1]

이러한 제품정보를 제공하려면 최종 제조사가 공급망의 앞선 단계에서 만들어진 근거를 확보하고 관리해야 한다.
배터리 규정 제39조는 셀·모듈 공급자에게 제조자의 규정 준수에 필요한 정보와 문서를 제공하도록 하고, 제77조 제4항은 책임 사업자에게 여권 정보의 정확성·완전성·최신성을 확보할 책임을 부여한다. [1]
그런데 원료가 여러 제품으로 나뉘거나 다른 원료와 함께 가공되면, 공급자가 제공한 근거를 각 제품에 어떻게 배분하고 반영했는지도 확인해야 한다.
증빙이 유효하게 발급됐다는 사실과 그 증빙을 이용한 후속 계산이 타당하다는 사실은 별개의 검증 대상이다.
본 연구는 이 중 생산·이전 과정의 회계 근거가 최종 제품의 주장에 사용되기까지의 관계를 검증하고 관리하는 문제를 다룬다. [2]

근거의 검증에는 정보 보호도 함께 요구된다.
배터리 규정은 공개 정보와 접근이 제한된 정보를 구분하며 데이터의 무결성·보안·프라이버시를 요구한다.
Babel 등도 공급망 정보가 통제되지 않은 채 확산될 수 있다는 우려를 DPP의 정보 공유를 어렵게 하는 요인으로 지적한다. [1, 2]
이에 본 연구는 공개할 제품정보와 이를 뒷받침하는 상세 거래자료를 구분한다.
기업은 민감한 회계 수치와 실제 조달·생산 연결을 일반 검증자에게 직접 제공하지 않으면서, 그 자료를 사용한 거래와 최종 주장 발급이 정해진 규칙을 따랐음을 검증받는다.
정상 거래를 검증하는 권한과 상세 이력을 열람하는 권한을 나누는 운영이다.


\subsection{정상 검증을 보완하는 사후 대응}
이 운영에서 원천 증빙의 타당성은 외부의 공급·검사·감사 절차에서 판단하고, 거래 검증은 수락된 자료의 계산·배분·이전이 정해진 규칙을 지켰는지 확인한다.
공급자의 증빙을 바탕으로 생산과 이전이 진행된 뒤, 추가 검사에서 그 증빙의 오류가 발견되는 상황을 생각할 수 있다.
자료에 오류가 있더라도 그 자료를 사용한 계산은 규칙에 맞을 수 있다.
따라서 정상 거래의 검증과 함께, 해당 자료를 계속 근거로 인정할 것인지 재평가하고 그 판단을 이후 사용에 반영하는 절차가 필요하다.

권한 있는 감사자가 특정 근거의 추가 사용을 중단하기로 결정했다면, 그 결정은 현재 남아 있는 파생 기록의 사용에도 적용돼야 한다.
이미 여러 기업의 생산·이전을 거친 경우 최초 기록만 표시해서는 이후 거래나 새로운 주장 발급에 사용되는 기록을 제어할 수 없기 때문이다.
이때 평소 공개하지 않은 거래 연결을 조사하여 관련 기록을 찾고, 그 기록을 사용하는 다음 거래가 거부되도록 해야 한다.
감사자는 근거의 재평가와 중단 결정을 맡고, 시스템은 실제 수락 이력의 조사와 승인된 사용 제한을 지원한다.
본 연구는 이 사후 대응을 정상 검증을 보완하는 안전장치로 삼는다.

따라서 본 연구의 질문은 다음과 같다.
\textbf{민감한 공급망 정보를 일반 검증자에게 직접 공개하지 않으면서 회계 근거를 최종 주장에 연결하고, 그 근거에 대한 사후 판단을 후속 사용에 반영하려면 어떤 기록·검증·권한 구조가 필요한가?}


\subsection{선행연구와 본 연구의 과제}
기존 연구는 비공개 검증과 공급망 정보 관리에 필요한 기반을 제공한다.
Babel 등은 생산 단계의 자격증명을 연결하고 원료 지분의 이전과 소비를 관리하는 DPP 구성을 제시했다.
Udokwu와 Craß는 자격증명으로 확인한 입력을 이용해 지속가능성 주장을 증명·검증하는 체계를 다룬다.
Sahai 등은 분할·합류·가공을 포함하는 공급망 문서 관계와 추적을 제공하며, 입력 문서 참조를 원장에 기록한다.
비공개 자산 거래 프로토콜인 CAP(Configurable Asset Privacy)는 정책에 따른 열람·동결과 실제 거래에 결합된 감사 자료를 지원한다. [2, 3, 4, 5]

본 연구가 구체화하는 대상은 이러한 기법을 사용하여 \textbf{DPP의 발급 근거와 사후 사용 통제를 같은 공급망 이력에 연결하는 시스템}이다.
최종 제품에 대한 검증은 Claim이 유효한 회계 상태와 발급 조건에 근거했는지를 확인해야 한다.
반대로 문제 출처에서 시작하는 사후 조사는 그 근거가 어느 후속 기록으로 이어졌으며 지금 무엇의 사용을 제한해야 하는지 확인해야 한다.
두 업무가 같은 이력을 다루려면 거래에서 검증한 근거, 최종 발급에 사용한 상태, 감사에서 복원한 관계와 원장의 사용 검사 대상이 서로 대응해야 한다.
기존 비공개 거래 기법을 공급망 응용으로 확장할 때에도 이 대응과 그에 따른 권한·처리 규칙을 정해야 한다.


\subsection{zkDPP의 설계와 기여}
본 논문은 이 과제를 위한 시스템인 zkDPP를 제안한다.
zkDPP는 생산·분할·가공·기업 간 이전을 공유 원장의 상태 변화로 표현한다.
참여자는 실제 입력과 출력이 회계·소유·권한 규칙을 만족함을 영지식증명으로 보이고, 원장은 이 증명과 입력 기록의 현재 사용 가능 상태를 함께 검사한다.
최종 발급에서는 회계 상태가 발급 정책을 충족하는지 확인한 뒤 해당 기록을 소비하고, 지정된 제품정보와 정책에 연결된 공개 주장 참조인 DPP Claim을 등록한다.
외부 검증자는 상세 거래자료를 열람하지 않고도 이 참조가 해당 제품정보 및 원장의 발급 이력에 결합됐는지 확인한다.

사후 조사를 위해 각 거래는 실제 입력 연결과 이후 소비를 확인하는 데 필요한 정보를 암호화해 남긴다.
거래 증명은 이 감사 자료가 해당 거래의 실제 입력·출력과 일치함도 검사한다.
이로써 감사자는 거래 당시 검증된 관계를 따라 이력을 조사할 수 있다.
감사 데이터의 복호화 키는 $N$명의 위원 중 $K$명 이상의 협력이 필요한 임계 암호화로 감싸 보호하고, 외부 승인 절차를 거쳐 복구·전달한다.
$K$가 2 이상이고 초기 키 사본이 올바르게 관리되는 조건에서 위원 한 명의 키 조각만으로는 감사 자료를 열 수 없다.
승인된 감사자는 전달받은 키로 보호된 기록의 연결을 로컬에서 조사하며, 후속 조치의 판단과 집행을 담당한다.

사용 중단이 승인되면 감사자는 확정된 원장 시점에서 문제 출처의 후속 기록을 추적하여, 아직 소비되지 않은 기록 전체를 사용 제한 대상으로 정한다.
여기에는 아직 완료되지 않은 기업 간 이전도 포함한다.
그 기록이 동결된 동안에는 후속 가공·이전과 새로운 Claim 발급이 거부된다.
발급과 동결은 공유 원장의 승인 순서와 같은 사용 상태에 따라 처리된다.
동결이 먼저 확정되면 차단 유지 중 해당 기록의 발급을 거부하고, 발급이 먼저 확정되면 이를 완료된 발급 이력으로 남긴다.
따라서 여러 갈래로 나뉜 이력에서도 이미 발급한 Claim을 보고하는 일과, 아직 소비되지 않은 가지의 다음 사용을 제한하는 일을 구분할 수 있다.

차단 결과는 처음 정한 전체 대상을 보존한 채 하나의 공통 확정 시점에서 확인한다.
대상이 있는 경우 그 전체가 여전히 미소비이면서 차단 상태이고 처리 결과가 명확한 때에만 전체 완료를 보고한다.
대상이 먼저 소비됐거나 조회·제출 결과가 불명확하면 그 사유와 대상을 남기고, 차단할 대상이 없는 경우는 별도로 보고한다.
이를 통해 일부 동결 성공을 전체 사용 제한의 완료로 해석하는 일을 방지한다.

본 논문의 기여는 다음과 같다.
\begin{enumerate}
    \item \textbf{비공개 회계·Claim 발급·승인된 사후 대응을 연결하는 시스템 설계.}
          생산과 이전에 따른 근거의 변화를 검증하고 이를 최종 Claim에 연결하며, 감사키 접근과 후속 사용 제한까지 수행하는 프로토콜을 정의한다.
          같은 기록에 대한 발급과 동결, 완료된 발급과 남은 기록의 처리를 일관된 전이·권한 규칙으로 명세한다.
          이는 해당 운영을 구현할 때 필요한 증명, 원장 검사와 감사 절차의 관계를 제공한다.
    \item \textbf{시스템의 보장과 실패 조건에 대한 조건부 분석.}
          명시한 암호·원장·권한 가정 아래 정확성, 건전성, 프라이버시를 정의하고 증명한다.
          실제 생성된 기록을 기준으로 거래·Claim·감사 자료의 일치를 보이며, 이를 정확한 추적과 고정된 전체 대상의 사용 차단으로 연결한다.
          공통 확인 시점에 전체 대상이 미소비·차단 상태임을 확인했을 때 차단 유지 중 추가 사용이 거부됨을 보이고, 자료 실패와 소비 경합에 따른 미완료도 구분한다.
          프라이버시는 참여자 비밀과 감사·상태 정보의 공개를 제외한 정상 이력의 수동 외부 관찰을 대상으로 분석한다.
\end{enumerate}

